\documentclass[11pt]{article}
\usepackage[T1]{fontenc}
\usepackage[utf8]{inputenc}
\usepackage[english]{babel}
\usepackage[margin=1in]{geometry}
\usepackage{amsmath,amssymb}
\usepackage{enumitem}
\usepackage{microtype}
\usepackage[hidelinks]{hyperref}

\setlength{\parindent}{0pt}
\setlength{\parskip}{0.65em}
\setlist{leftmargin=*,topsep=0.3em,itemsep=0.35em}

\title{Technical Review of\\
\emph{Early-Stopping Activation Steering for Factual Preference Alignment in Vietnamese Medical Language Models}}
\author{}
\date{September 3, 2026}

\begin{document}
\maketitle

\section*{Overall assessment}

The manuscript reports a potentially useful inference-time intervention and supplies paired outcome counts for several central comparisons.  The strongest internally supported completed-answer result is the Hard Cutoff improvement under \texttt{max\_new\_tokens=800}: the displayed discordant counts reproduce the reported exact McNemar result, and the Holm-adjusted conclusion is coherent.  The bounded Linear Decay comparison also has a reproducible exact McNemar result, but the ablation table indicates that essentially all 200-token outputs are truncated rather than naturally completed.

Major revision is still needed.  The activation-norm paragraph contains definite arithmetic errors and is inconsistent with the claimed fixed teacher-forcing design.  The post-prefill hook description implies an off-by-one ambiguity in which output tokens are steered; the negative control has a double-sign ambiguity; and the stated Tesla T4/\texttt{bfloat16} environment is not compatible with native CUDA BF16 execution.  The control suite does not yet isolate a truthfulness direction, and the main factual claims remain stronger than the BERTScore reference-preference endpoint.  Natural-completion evidence favors Hard Cutoff, not the specifically promoted Linear Decay schedule.  Statistical interval labels, RAG superiority language, and multiple bibliography records also require correction.

\section*{Major technical findings}

\subsection*{1. The activation-norm arithmetic is wrong, and the experiment design is internally inconsistent}

\textbf{Classification: definite errors and implementation inconsistency.}

\textbf{Location.} Abstract; contribution 3; Empirical Activation-Norm Dynamics; Discussion.

\textbf{Problem.} The reported values imply
\begin{align*}
56.28-53.64 &= +2.64 \quad (t=10,\ \text{Continuous}),\\
56.34-53.39 &= +2.95 \quad (t=50,\ \text{Continuous}),\\
52.95-53.64 &= -0.69 \quad (t=10,\ \text{Linear Decay}),\\
53.39-53.39 &= 0 \quad (t=50,\ \text{Linear Decay}),\\
52.75-53.39 &= -0.64 \quad (t=50,\ \text{Hard Cutoff}).
\end{align*}
The prose instead states a ``sustained $+1.97$'' Continuous elevation and a Linear Decay ``norm undershoot of $-3.01$ units relative to baseline $53.39$.''  Neither statement follows from the displayed numbers.

There is also a design contradiction.  Contribution 3 calls this a fixed teacher-forcing experiment.  If the same tokens are teacher-forced and the hook is the identity for $t>K$, both Linear Decay and Hard Cutoff must have the same measured Layer~8 state as baseline at $t=50$.  Linear Decay does, but Hard Cutoff does not.  If the values instead come from free-running generations, post-$K$ differences reflect changed token histories and cannot be interpreted as direct relaxation of the hook.  Two time points also do not establish an ``abrupt transition'' or an equilibrium.

\textbf{Why it matters.} Norm mitigation is the proposed mechanism and a named contribution.  The current arithmetic and incompatible design descriptions prevent verification and do not establish a causal link between the hook, cumulative activation growth, or output degeneration.

\textbf{Suggested fix.} Recompute every difference and replace all norm values from one documented experiment.  State whether tokens are fixed or free-running; report pre-hook and post-hook norms, $\|\hat h-h\|_2$, $\langle h,v_{\text{steer}}\rangle$, uncertainty across prompts, and the full trajectory.  Under teacher forcing, describe equality after $K$ as an expected consequence of $\alpha(t)=0$, not independent evidence of behavioral mitigation.

\subsection*{2. Token timing, layer naming, and the negative-control sign are implementation-critical ambiguities}

\textbf{Classification: likely off-by-one implementation error and definite notation/sign ambiguity.}

\textbf{Location.} Early-Stopping Steering Hook; Experimental Setup; abstract and contribution list; Directional Controls.

\textbf{Problem.} The manuscript states that prompt prefill ($t=0$) is unsteered and that the first post-prefill decoding call is $t=1$.  With standard cached autoregressive generation, prefill logits sample the first output token.  The $t=1$ call consumes that token and produces logits for the second output token.  The described hook therefore affects predicted output tokens 2 through $K+1$, not the ``first $K$ generated tokens,'' unless a different custom decoding loop is used.

Layer terminology drifts between ``Layer~8,'' $h_8$, and \texttt{model.model.layers[8]}, which the paper correctly identifies as the ninth decoder block.  Use of a human layer number and a zero-based module index as if they were identical is error-prone.  The negative control has a related sign ambiguity: it is denoted $-v_{\text{steer}}$ while also assigning $\alpha_0=-18.0$.  Applying both signs in $h+\alpha v$ would produce positive steering.

The extraction state is the last token of $x_i\oplus y_i$, but the chat template, answer-boundary token, EOS inclusion, response-length handling, and exact equivalence of extraction and injection sites are not specified.  The dose formulas $D_{\text{cutoff}}=K|\alpha_0|$ and $D_{\text{decay}}=(K+1)|\alpha_0|/2$ also silently assume $T\ge K$; for shorter generations the upper limit is $\min(T,K)$.

\textbf{Why it matters.} These choices determine which logits are changed, which block is intervened on, and even the direction of the principal control.  They can materially alter every reported output.

\textbf{Suggested fix.} Supply executable pseudocode for prefill and each cached decoding call, explicitly naming the input token and output token whose logits are affected.  Use ``decoder block at zero-based index 8 (ninth block)'' consistently.  Implement negative steering as either $(v,\alpha)=\left(v_{\text{steer}},-18\right)$ or $\left(-v_{\text{steer}},+18\right)$, not both.  Specify tokenization, chat template, EOS handling, tensor slice, tuple reconstruction, dtype/device casts, and the $T<K$ schedule cases.

\subsection*{3. The stated Tesla T4 and \texttt{bfloat16} environment is not self-consistent}

\textbf{Classification: externally verified likely implementation/reporting error.}

\textbf{Location.} Model Checkpoint \& Precision; Reproducibility Environment; latency claims.

\textbf{Problem.} The manuscript reports NF4 with \texttt{bfloat16} compute on NVIDIA Tesla T4 GPUs.  NVIDIA identifies Turing devices as compute capability 7.5, while the CUDA programming guide requires compute capability 8.0 or higher for native \texttt{bfloat16}.  NVIDIA's T4 specifications list FP16/FP32 mixed precision, INT8, and INT4, but not BF16.  Thus the stated configuration cannot describe native CUDA BF16 execution on a T4.  The run may have used FP16, a fallback/conversion path, different hardware, or incorrectly reported metadata.

\textbf{Why it matters.} Compute dtype affects numerical behavior, memory, kernel selection, and latency.  This inconsistency directly undermines reproduction and the repeated systems-overhead comparisons.

\textbf{Suggested fix.} Report the exact output of \nolinkurl{torch.cuda.get_device_capability()} and \nolinkurl{torch.cuda.is_bf16_supported()}, together with model parameter dtypes, BitsAndBytes configuration, and autocast state from the run logs.  If the runs used T4, use and report FP16 unless a documented emulation/fallback path was actually used; if BF16 was used natively, correct the GPU model.  Re-run latency measurements after reconciling the environment.

\subsection*{4. The directional controls do not yet establish truthfulness specificity}

\textbf{Classification: definite control-design error and likely validity problem.}

\textbf{Location.} Contrastive Vector Estimation \& Directional Controls; Reproducibility Environment; Table~\texttt{tab:baselines}; directional-specificity claims.

\textbf{Problem.} Pair shuffling is invariant for this mean-difference estimator:
\[
\frac{1}{N}\sum_i\left(h_i^+-h_{\pi(i)}^-\right)=\mu^+-\mu^-.
\]
Its reported cosine of 1.0000 is therefore mathematically inevitable, so it is not a negative control; it is also absent from the results table.  Label shuffling is unclear: $N_{\text{flipped}}=185$ is not reconciled with the stated 10,290 training pairs.  If only 185 of 10,290 labels are flipped, most of the learned signal remains; if the vector used a smaller construction subset, that subset is undocumented.

The random-vector definitions are incomplete.  $z_r$ is undefined in $v_{\text{cov}}^{(r)}=Lz_r/\|Lz_r\|_2$.  The manuscript writes normalized $v_{\text{rand}}$ as Gaussian, although normalizing a Gaussian draw produces a unit-sphere direction; it should define $z_r\sim\mathcal N(0,I)$ followed by $v_r=z_r/\|z_r\|_2$.  The schedule and coefficient applied to the shuffled, covariance-matched, and isotropic controls are not stated explicitly.

The isotropic distribution is also a weak placebo for the claimed inference: its mean RefPref is 55.82\%, far below the 73.00\% unsteered baseline, so large arbitrary perturbations are generally harmful.  Showing that the learned vector is less harmful than these draws is not by itself evidence of truth-specific improvement.  The covariance-matched mean is 74.81\%, only 2.39~pp below the learned vector, but no vector-level range, exceedance count, paired improvement distribution, or pre-specified Monte Carlo test is reported.

\textbf{Why it matters.} Directional specificity is a central contribution.  The current controls conflate truth-related structure, answer-style/length artifacts, and generic damage caused by a large residual perturbation.

\textbf{Suggested fix.} Replace pair shuffling with full balanced label permutations or random signs on pairwise differences, using the complete documented construction set.  Match every control on layer, schedule, norm, coefficient, and evaluation prompts.  Define all random variables, report vector-level paired changes relative to baseline, and give the number of covariance-matched and label-permuted controls that equal or exceed $+v_{\text{steer}}$.  Add extraction controls matched for answer length, final token, and surface style.

\subsection*{5. Benchmark documentation and the endpoint do not support clinical factual-accuracy language}

\textbf{Classification: unsupported claim and likely dataset-validity issue.}

\textbf{Location.} Benchmark Construction; BERTScore Reference-Preference Criterion; abstract, Results, table captions, and conclusion.

\textbf{Problem.} The manuscript does not document the qualifications or number of experts, source-to-question procedure, synthetic-negative generator and prompts, factual review, agreement, adjudication, or correction process.  The category counts 165, 160, and 175 sum to 500, so they describe the evaluation subset, but the construction paragraph introduces them as if they characterize all 14,700 records.  The relationship between the 2,205-example test partition and the selected 500 questions is not defined, nor is ``question-disjoint'' specified by drug, entity, or source passage.  Overlap at those levels could let source- or style-specific structure enter both vector estimation and testing.

RefPref only tests whether an output is semantically closer to one reference than to one synthetic counter-answer.  It does not verify dosage arithmetic, contraindications, drug interactions, or new unsupported statements, and it can respond to length and style.  The paper acknowledges this limitation but repeatedly switches the endpoint's descriptor from ``reference preference'' to ``factual accuracy,'' ``factual alignment,'' or ``clinical safety.''  Tie counts and complete BERTScore settings, including IDF and baseline-rescaling choices, are not reported.

\textbf{Why it matters.} In a medical setting, a semantic preference proxy is not a clinical correctness or safety endpoint.  Synthetic artifacts or source leakage could explain part of the measured effect.

\textbf{Suggested fix.} Document the complete benchmark pipeline, expert validation, category counts for every split, the 500-example sampling rule, and overlap checks.  Group splits by source passage and drug/entity and release immutable split IDs.  Add blinded clinician adjudication and structured scoring for dose, interaction, contraindication, and unsupported-claim errors.  Report ties and all BERTScore options, and use ``RefPref rate'' consistently for the current metric.

\subsection*{6. Schedule conclusions are conflated with a generation-cap effect}

\textbf{Classification: claim/evidence mismatch and terminology drift.}

\textbf{Location.} Abstract; both primary evaluations; Controlled Factorial Ablation; conclusion.

\textbf{Problem.} In natural completion, Hard Cutoff is the only schedule with a Holm-adjusted significant gain; Linear Decay gains 2.40~pp with $p_{\text{adj}}=0.1189$.  In the 200-token experiment, Linear Decay is emphasized, but the ablation table reports 0.0\% EOS for that condition and at most 0.2\% for any condition.  These are almost entirely truncated prefixes.  The unsteered baseline changes from 73.00\% at 200 tokens to 65.40\% under natural completion, demonstrating strong cap sensitivity.

The factorial table does not report RefPref at all, even though RefPref is the primary endpoint, and it supplies no paired uncertainty for schedule differences.  Thus it cannot establish which temporal schedule is best for the claimed factual-preference effect.  The term ``early-stopping steering'' alternates between Linear Decay alone and an umbrella covering both Linear Decay and Hard Cutoff.

\textbf{Why it matters.} A statistically significant result on truncated prefixes cannot be assumed to transfer to completed answers, and the completed-answer experiment currently supports a different schedule from the one highlighted in the title narrative.

\textbf{Suggested fix.} Pre-specify a primary schedule and endpoint.  Call the 200-token analysis a truncated-prefix stress test, report output-length distributions, and evaluate all schedules with RefPref and paired uncertainty at both caps.  State plainly that current natural-completion evidence favors Hard Cutoff and treats Linear Decay as exploratory.  Use schedule-specific terminology throughout.

\subsection*{7. The claimed repetition mechanism is contradicted by the reported results, and ``Experiment 10'' is unsupported}

\textbf{Classification: unsupported causal claim and unsupported result.}

\textbf{Location.} Introduction; natural-completion results; Empirical Activation-Norm Dynamics; Multi-Metric Pareto Trade-off \& Deployment Mitigation.

\textbf{Problem.} The Introduction says continuous steering's norm offset is ``manifested as repetitive text loops and degraded syntax.''  No syntax measure or loop analysis is reported.  The only completed-answer repetition metric points in the opposite direction: Rep-4 is 4.38\% for Continuous Steering, compared with 5.35\% for Hard Cutoff and 5.37\% for Linear Decay.  EOS incidence also cannot prove the absence of repetitive loops.

The Discussion introduces ``Experiment 10,'' a repetition penalty $\theta_{\text{rep}}=1.15$, Rep-4 of 3.76\%, and latencies of 78.74~s and 45.05~s.  No such experiment, method, table, test-set definition, output-length result, uncertainty, or RefPref result appears elsewhere.  Therefore the claims that the trade-off is ``resolved completely'' and that the $+5.20$~pp gain is preserved are unsupported.  The paragraph subsequently repeats the unpenalized trade-off and calls the same penalty only a ``viable path,'' which is internally inconsistent.

\textbf{Why it matters.} The manuscript's novelty is motivated as preventing degeneration caused by continuous steering, but the presented behavioral evidence neither demonstrates that mechanism nor validates the proposed remedy.

\textbf{Suggested fix.} Remove the causal and Experiment~10 claims, or add a complete matched experiment for every schedule with RefPref, reference BERTScore, Rep-4, syntax/degeneration measures, EOS, generated-token counts, and synchronized latency.  Until then, limit the conclusion to a post-hook norm offset whose relationship to output quality remains unresolved.

\subsection*{8. Confidence-interval labels and endpoint descriptions need reconciliation}

\textbf{Classification: likely statistical-reporting error; exact tests themselves are internally correct.}

\textbf{Location.} Contribution 5; Tables~\texttt{tab:mcnemar} and \texttt{tab:mcnemar\_bounded}; RAG comparison; Discussion.

\textbf{Problem.} The displayed exact two-sided McNemar values for $(b,c)=(16,37)$, $(16,42)$, $(18,34)$, $(19,31)$, and $(22,27)$ are reproducible.  However, the reported intervals $[+1.37,+7.03]$~pp and $[+2.25,+8.15]$~pp numerically equal the paired normal intervals
\[
\widehat\Delta\ \pm\ 1.96\sqrt{\frac{b+c-(c-b)^2/N}{N^2}},
\]
while the natural-completion table labels its intervals as a $B=10{,}000$ percentile bootstrap with seed 42.  A direct paired percentile resampling of the displayed four-cell counts does not reproduce the stated endpoints under common seeded implementations.  The per-example data and code needed to resolve this are absent.

The contribution list also says that Wilcoxon testing ``confirm[s] reference-preference gains,'' but the Wilcoxon test concerns continuous per-example BERTScore F1, not the binary RefPref endpoint.  In the Hybrid RAG matrix, the prose calls the discordant cells $b=22,c=27$, whereas the usual row-by-column convention for the displayed matrix would label the upper-right cell 27 as $b$ and the lower-left cell 22 as $c$.  The two-sided $p$-value is unchanged, but the notation drifts across tables.

\textbf{Why it matters.} Statistical validation is a named contribution.  Mislabeling the interval method or conflating binary and continuous endpoints obscures what has actually been tested.

\textbf{Suggested fix.} Release the 500 paired outcomes, per-example BERTScore values, and exact analysis code.  Name each interval method correctly and report a paired effect interval for every principal comparison.  Describe Wilcoxon as a separate BERTScore-distribution analysis, including zero-difference handling.  Define discordant-cell orientation once or avoid $b,c$ labels when the table orientation changes.

\subsection*{9. The Hybrid RAG comparison does not support an ``outperforms'' claim}

\textbf{Classification: unsupported comparative claim and reproducibility ambiguity.}

\textbf{Location.} Placebo/RAG Results; Tables~\texttt{tab:baselines} and \texttt{tab:mcnemar\_rag}; abstract and conclusion.

\textbf{Problem.} The manuscript first says Linear Decay ``outperforms Hybrid RAG'' by 1.00~pp, then reports an exact paired $p=0.5682$ and correctly calls the RefPref rates statistically indistinguishable.  Hybrid RAG also has higher absolute reference BERTScore F1 (0.7190 versus 0.7150).  Steering offers lower reported latency and context use, but it does not dominate Hybrid RAG across the reported metrics.

The retrieval comparison is not fully reproducible: passage construction, retrieval top-$k$, dense pooling and score normalization, model revision, RRF inputs, context-selection and truncation rules, prompts, and the \texttt{pyvi} version are missing.  It is also unclear whether latency includes retrieval, whether CUDA synchronization and warm-up were used, and whether output lengths were matched.  Repeating exactly $7.32\pm0.84$~s for every non-RAG condition does not by itself establish zero hook overhead.

\textbf{Why it matters.} The strongest non-oracle RAG comparator is statistically tied on RefPref and better on absolute BERTScore.  Timing and context trade-offs are useful, but they require a controlled systems measurement rather than a generic superiority claim.

\textbf{Suggested fix.} Replace ``outperforms'' with ``numerically exceeds'' and present Hybrid RAG as a quality--latency--context trade-off.  Fully specify retrieval and timing protocols, report output-token distributions, decompose retrieval/prefill/decoding latency, and give uncertainty for paired latency differences.  Keep the BM25 comparison explicitly limited to that weaker lexical configuration.

\subsection*{10. The layer-sweep interpretation exceeds the evidence}

\textbf{Classification: unsupported inference.}

\textbf{Location.} Layer-Wise Subspace Probing.

\textbf{Problem.} A BERTScore range of 0.7040--0.7140 over layers 4--15 on $N_{\text{eval}}=100$, without the individual values, paired uncertainty, or repeated selection, does not show that truthfulness representations are distributed across layers and does not ``confirm framework robustness.''  Similar downstream values may arise from residual propagation, correlated layers, or sampling noise.

\textbf{Why it matters.} The paragraph converts an exploratory hyperparameter sweep into a mechanistic localization and robustness conclusion.

\textbf{Suggested fix.} Report the complete sweep, paired intervals, search criterion, and all other selected hyperparameters.  Limit the conclusion to the observation that several tested intervention layers produced similar validation scores unless a dedicated localization experiment is added.

\section*{Reference integrity and meaningful editorial findings}

Several bibliography records do not match the externally verifiable primary-source metadata and should be re-exported rather than edited piecemeal:

\begin{enumerate}
 \item \textbf{Ref.~\texttt{ref1}:} the author list does not match the official record for arXiv:2310.01405, \emph{Representation Engineering: A Top-Down Approach to AI Transparency}.
 \item \textbf{Ref.~\texttt{ref2}:} the official NeurIPS record lists Kenneth Li, Oam Patel, Fernanda Vi\'egas, Hanspeter Pfister, and Martin Wattenberg.  The manuscript omits Pfister and incorrectly includes N.~Nanda.
 \item \textbf{Ref.~\texttt{ref4}:} arXiv:2409.12186 is the \emph{Qwen2.5-Coder Technical Report}, not the cited \emph{Qwen2.5 Technical Report}.  The latter is arXiv:2412.15115.  The manuscript's 28-block and 3584-dimensional configuration is consistent with the official model configuration, but the citation identifier is wrong.
 \item \textbf{Ref.~\texttt{ref8}:} the current arXiv:2308.10248 metadata includes Juan J.~Vazquez and differs from the manuscript's author list/order; update it from the primary record.
 \item \textbf{Ref.~\texttt{ref9}:} the ACL 2024 paper's authors are Nina Rimsky, Nick Gabrieli, Julian Schulz, Meg Tong, Evan Hubinger, and Alexander Turner, with pages 15504--15522.  The manuscript lists different coauthors and omits the page range.
\end{enumerate}

The exact software/environment sentence is duplicated in consecutive paragraphs in Experimental Setup.  Delete one copy and retain the random-seed information in the remaining paragraph.  The multi-metric Discussion paragraph likewise repeats the natural-completion trade-off; consolidate it after resolving the unsupported Experiment~10 material.

\section*{Proofing sweep}

The bundled mechanical scan was run on both the compiled PDF and the TeX source.  Its PDF double-period hit is the legitimate ellipsis in $\{1,\ldots,100\}$, and its semicolon-capitalization hit is the proper metric name ``Recall@3''; neither is an error.  No inverse-trigonometric or coordinate-transform expression appears, so no branch or quadrant ambiguity was found.

High-confidence local corrections are:

\begin{itemize}
 \item \textbf{Natural-completion Results:} 5.35\% versus 4.10\% is a 30.5\% relative increase, but 5.37\% versus 4.10\% is approximately 31.0\%.  Name the schedule or report both values.
 \item \textbf{RAG Results:} the Hybrid RRF sentence uses ``achieving'' twice in one clause.  Replace the second occurrence with ``and yields.''
 \item \textbf{Units:} write latency consistently as $7.32$~s and $+7.13$~s rather than \texttt{7.32s} and \texttt{+7.13s}.
 \item \textbf{Control organization:} ``we evaluate two controls'' is followed by a second numbered item containing five conditions and repeating Negative Steering.  Replace it with one flat list whose entries match the reported table.
 \item \textbf{Typesetting:} the manuscript compiles, but the log reports a 22.16~pt overfull box around the Hybrid contingency table.  Reduce column spacing, abbreviate headings, or resize that table without making it unreadably small.
\end{itemize}

\section*{Highest-priority revision sequence}

\begin{enumerate}
 \item Correct and redesign the activation-norm analysis, then remove unsupported degeneration and Experiment~10 claims.
 \item Resolve token timing, layer indexing, negative-control sign, hardware/dtype reporting, and all control-vector definitions in executable pseudocode and run logs.
 \item Replace invariant/partial shuffles with valid matched null controls and report vector-level paired effects.
 \item Document the benchmark and source-level split, and add clinician-adjudicated factual evaluation.
 \item Pre-specify the primary schedule and separate completed-answer evidence from truncated-prefix evidence.
 \item Correct the interval methodology, Hybrid RAG language, timing protocol, and bibliography metadata.
\end{enumerate}

\section*{Items requiring external verification or unavailable artifacts}

The following cannot be established from the manuscript alone and should be verified before publication:

\begin{itemize}
 \item expert qualifications, factual adjudication, public availability, and redistribution permission for ViHaluEval-Medical and the formulary-derived passages;
 \item whether Ref.~\texttt{ref3} supports the specific claims about drug-interaction, dosage, and contraindication errors, rather than only broader clinical-knowledge limitations;
 \item whether Ref.~\texttt{ref11} supports the asserted internal-activation explanation for ignoring retrieved evidence;
 \item run logs showing the actual GPU, CUDA capability, compute dtype, fallback behavior, model revision, and loaded quantization configuration;
 \item source code, immutable split IDs, retrieval index, per-query outputs, and statistical scripts needed to reproduce the norm, latency, RAG, random-control, Wilcoxon, and bootstrap claims.
\end{itemize}

\section*{Checks that were internally consistent}

The 70/15/15 split counts sum to 14,700; the three evaluation categories sum to 500; category-level preferred counts sum to the stated totals; $D_{\text{decay}}=(K+1)|\alpha_0|/2$ and $\alpha_{\text{match}}=0.0425\alpha_0$ are algebraically correct for $K=16$, $T=200$, and $T\ge K$; and the displayed exact McNemar $p$-values agree with their discordant counts.  The official Qwen2.5-7B-Instruct configuration also supports the stated 28 decoder blocks and hidden size 3584.  No dimensional error was found in the steering equations themselves, and no inverse-trigonometric branch issue is present.  These clean checks do not resolve the sign, timing, hardware, control, and endpoint ambiguities above.

\end{document}